% !PRINTED 260
\weberpage{260}
%<<CONT
und $\delta$ und $\beta$ lassen sich so bestimmen, daß

\weq{7}{\delta(\alpha+\beta) - \beta(\gamma+\delta) = \alpha\delta - \beta\gamma = 1.}

Aus den Gleichungen (6) für $a$ und $c-\partial$ findet sich dann

\weq{8}{c'+a' = a\delta - (c-\partial)\beta.}

$\delta$ und $\beta$ können, da $\alpha+\beta$ und $\gamma+\delta$ beide ungerade sind,
nach (7) nicht beide ungerade sein; folglich fällt $c'$ nach (8)
gerade aus. Ersetzt man, was nach (7) gestattet ist, $\delta$, $\beta$ durch
$\delta+h(\gamma+\delta)$, $\beta+h(\alpha+\beta)$, für ein beliebiges $h$, so ändert
sich $c'$ nach (6) und (8) um $2ha'$, und über $h$ kann so verfügt
werden, daß $c'$ durch 16 oder 48 teilbar wird.

Es ist dann nach (6) $\alpha+\beta+\gamma-\delta$ gerade und daher
die Formel (12), \S\,40 anzuwenden. Darin ist nun zu berücksichtigen


\[
\begin{aligned}
a'(\alpha-\beta) &\equiv a, & \partial(\alpha-\beta) &\equiv \partial' \\
\partial'(\alpha+\beta) &\equiv a, & a'(\gamma-\delta) &\equiv -\partial
\end{aligned} \quad (\mathrm{mod}\ 16),
\]

woraus folgt:

\[
n(\alpha-\beta)(\alpha+\beta+\gamma-\delta) \equiv a^2 - \partial'^2 \ (\mathrm{mod}\ 16),
\]

also

\[
\left(\frac{2}{\alpha-\beta}\right) e^{-\frac{3\pi i}{8}(\alpha-\beta)(\alpha+\beta+\gamma-\delta)} = \left(\frac{2}{a}\right)\left(\frac{2}{a'}\right)\left(\frac{2}{a}\right)\left(\frac{2}{\partial'}\right) = \left(\frac{2}{n}\right)
\]

und

\[
\varrho = e^{-\frac{2\pi i}{3}\left[\alpha(\gamma-\beta)-(\alpha^2-1)\beta\delta\right]}.
\]

Ist $n$ nicht durch 3 teilbar, so ist

\[
\begin{aligned}
2\alpha &\equiv a(a'+\partial'), & 2\delta &\equiv \partial(a'+\partial'), \\
2\beta &\equiv a(\partial'-a'), & 2\gamma &\equiv \partial(\partial'-a')
\end{aligned} \quad (\mathrm{mod}\ 3),
\]

also entweder $\alpha$ und $\delta$ oder $\beta$ und $\gamma$ durch 3 teilbar und also
$\varrho = 1$.

Hieraus ergeben sich folgende zusammengehörige Vertauschungen
 (wobei die Änderung von $a, c, \partial$ in $a', c', \partial'$ durch (6)
bestimmt ist):

\weq{9}{\begin{array}{ccc}
\omega, & u, & v \\[4pt]
\dfrac{\omega-1}{\omega+1}, & \dfrac{\sqrt{2}}{u}, & \varrho\left(\dfrac{2}{n}\right)\dfrac{\sqrt{2}}{v}.
\end{array}}

Wir wenden die Vertauschungen (4), (5), (9) auf folgende
Funktionen an:
%CONT>>
